% GENERATED FILE. Edit canonical lessons, then run npm run generate:books.
% canonical-source-hash: fnv1a64:4c97c138ee9f8822
% canonical-lessons: BN-C123-beral, BN-C123-konui, BN-C123-buk, BN-C123-komor, BN-C123-gorali

\chapter{Cat, Elbow, Chest, Waist, Ankle}
\label{ch:bn-cat-elbow-chest-waist-ankle}
\hlchaptermodality{123}

\begin{chapteropening}
\textbf{By the end of this chapter:} \emph{I can say a cat, an elbow, the chest, the waist and an ankle.}

Complete the last lesson of chapter 123: I can say a cat, an elbow, the chest, the waist and an ankle.
\end{chapteropening}

\section[\texorpdfstring{\bn{বেড়াল} (beṛāl)}{বেড়াল (beṛāl)}]{\bn{বেড়াল} (beṛāl) — a cat}

\label{lesson:BN-C123-beral}



\begin{quote}
Before the new one: say the Bengali for a gecko, then the Bengali for a turtle.
\end{quote}

\subsection*{You'll want to know: \bn{বেড়াল}}
\textbf{\bn{বেড়াল}} — \emph{beṛāl} — \textquotedblleft{}a cat\textquotedblright{}.

\begin{cousinweb}
\textbf{\bn{বেড়াল}} (\emph{beṛāl}) is a tadbhava, a word inherited from Sanskrit and reshaped by sound change. It comes from \emph{biḍāla}, \textquotedblleft{}cat\textquotedblright{} (CDIAL 9237). CDIAL writes the Bengali forms \emph{birāl} and \emph{bilāl}.
\end{cousinweb}

\begin{grammarlens}[title={What you've built}]
One more word for this chapter.
\end{grammarlens}

\subsection*{Your turn}
\begin{itemize}
\raggedright
  \item \emph{Say it:} \emph{beṛāl}
  \item \emph{Say it:} \emph{beṛāl}, once more
  \item \emph{Say it:} say \emph{beṛāl}
  \item \emph{Recall:} say the Bengali for a gecko, then the Bengali for a turtle, then say \emph{beṛāl} again
\end{itemize}

\begin{tcolorbox}[breakable,colback=teal!4,colframe=teal!35!black,title={Before you move on}]
What does \bn{বেড়াল} mean? (\textquotedblleft{}A cat\textquotedblright{}.) Say it once more.
\end{tcolorbox}
\section[\texorpdfstring{\bn{কনুই} (kônui)}{কনুই (kônui)}]{\bn{কনুই} (kônui) — an elbow}

\label{lesson:BN-C123-konui}



\begin{quote}
Before the new one: say the Bengali for a turtle, then the Bengali for a cat.
\end{quote}

\subsection*{You'll want to know: \bn{কনুই}}
\textbf{\bn{কনুই}} — \emph{kônui} — \textquotedblleft{}an elbow\textquotedblright{}.

\begin{cousinweb}
\textbf{\bn{কনুই}} (\emph{kônui}) is a tadbhava, a word inherited from Sanskrit and reshaped by sound change. It comes from \emph{kaphoṇi}, \textquotedblleft{}elbow\textquotedblright{} (CDIAL 2757). CDIAL thinks the word more probably Munda than Dravidian.
\end{cousinweb}

\begin{grammarlens}[title={What you've built}]
One more word for this chapter.
\end{grammarlens}

\subsection*{Your turn}
\begin{itemize}
\raggedright
  \item \emph{Say it:} \emph{kônui}
  \item \emph{Say it:} \emph{kônui}, once more
  \item \emph{Say it:} say \emph{kônui}
  \item \emph{Recall:} say the Bengali for a turtle, then the Bengali for a cat, then say \emph{kônui} again
\end{itemize}

\begin{tcolorbox}[breakable,colback=teal!4,colframe=teal!35!black,title={Before you move on}]
What does \bn{কনুই} mean? (\textquotedblleft{}An elbow\textquotedblright{}.) Say it once more.
\end{tcolorbox}
\section[\texorpdfstring{\bn{বুক} (buk)}{বুক (buk)}]{\bn{বুক} (buk) — the chest}

\label{lesson:BN-C123-buk}



\begin{quote}
Before the new one: say the Bengali for a cat, then the Bengali for an elbow.
\end{quote}

\subsection*{You'll want to know: \bn{বুক}}
\textbf{\bn{বুক}} — \emph{buk} — \textquotedblleft{}the chest\textquotedblright{}.

\begin{cousinweb}
\textbf{\bn{বুক}} (\emph{buk}) is a tadbhava, a word inherited from Sanskrit and reshaped by sound change. It comes from \emph{bukka}, \textquotedblleft{}heart\textquotedblright{} (CDIAL 12064). CDIAL treats \emph{bukka} as a non-Aryan word that crossed with Sanskrit \emph{vṛkka}, \textquotedblleft{}kidneys\textquotedblright{}.
\end{cousinweb}

\begin{grammarlens}[title={What you've built}]
One more word for this chapter.
\end{grammarlens}

\subsection*{Your turn}
\begin{itemize}
\raggedright
  \item \emph{Say it:} \emph{buk}
  \item \emph{Say it:} \emph{buk}, once more
  \item \emph{Say it:} say \emph{buk}
  \item \emph{Recall:} say the Bengali for a cat, then the Bengali for an elbow, then say \emph{buk} again
\end{itemize}

\begin{tcolorbox}[breakable,colback=teal!4,colframe=teal!35!black,title={Before you move on}]
What does \bn{বুক} mean? (\textquotedblleft{}The chest\textquotedblright{}.) Say it once more.
\end{tcolorbox}
\section[\texorpdfstring{\bn{কোমর} (kômôr)}{কোমর (kômôr)}]{\bn{কোমর} (kômôr) — the waist}

\label{lesson:BN-C123-komor}



\begin{quote}
Before the new one: say the Bengali for an elbow, then the Bengali for the chest.
\end{quote}

\subsection*{You'll want to know: \bn{কোমর}}
\textbf{\bn{কোমর}} — \emph{kômôr} — \textquotedblleft{}the waist\textquotedblright{}.

\begin{cousinweb}
\textbf{\bn{কোমর}} (\emph{kômôr}) is a loan, usually traced to Persian \emph{kamar}, \textquotedblleft{}waist\textquotedblright{}.
\end{cousinweb}

\begin{grammarlens}[title={What you've built}]
One more word for this chapter.
\end{grammarlens}

\subsection*{Your turn}
\begin{itemize}
\raggedright
  \item \emph{Say it:} \emph{kômôr}
  \item \emph{Say it:} \emph{kômôr}, once more
  \item \emph{Say it:} say \emph{kômôr}
  \item \emph{Recall:} say the Bengali for an elbow, then the Bengali for the chest, then say \emph{kômôr} again
\end{itemize}

\begin{tcolorbox}[breakable,colback=teal!4,colframe=teal!35!black,title={Before you move on}]
What does \bn{কোমর} mean? (\textquotedblleft{}The waist\textquotedblright{}.) Say it once more.
\end{tcolorbox}
\section[\texorpdfstring{\bn{গোড়ালি} (goṛāli)}{গোড়ালি (goṛāli)}]{\bn{গোড়ালি} (goṛāli) — an ankle}

\label{lesson:BN-C123-gorali}



\begin{quote}
Before the new one: say the Bengali for the chest, then the Bengali for the waist.
\end{quote}

\subsection*{You'll want to know: \bn{গোড়ালি}}
\textbf{\bn{গোড়ালি}} — \emph{goṛāli} — \textquotedblleft{}an ankle\textquotedblright{}.

\begin{cousinweb}
\textbf{\bn{গোড়ালি}} (\emph{goṛāli}) is built on \emph{goṛ}, \textquotedblleft{}foot, leg\textquotedblright{}, which CDIAL lists under \emph{*gōḍḍa}, \textquotedblleft{}foot, leg, knee\textquotedblright{}, a form it reconstructs (CDIAL 4272); CDIAL does not list \textbf{\bn{গোড়ালি}} (\emph{goṛāli}) itself.
\end{cousinweb}

\begin{grammarlens}[title={What you've built}]
One more word for this chapter.
\end{grammarlens}

\subsection*{Your turn}
\begin{itemize}
\raggedright
  \item \emph{Say it:} \emph{goṛāli}
  \item \emph{Say it:} \emph{goṛāli}, once more
  \item \emph{Say it:} say \emph{goṛāli}
  \item \emph{Recall:} say the Bengali for the chest, then the Bengali for the waist, then say \emph{goṛāli} again
\end{itemize}

\begin{tcolorbox}[breakable,colback=teal!4,colframe=teal!35!black,title={Before you move on}]
What does \bn{গোড়ালি} mean? (\textquotedblleft{}An ankle\textquotedblright{}.) Say it once more.
\end{tcolorbox}
